\documentclass{article}
\usepackage{amsmath,amssymb,amsthm,mathtools,mathrsfs}
\begin{document}
\maketitle
\section{The original Laplacian control along the exact period deformation}
\label{sec:period-laplacian-control-bridge}

We calculate the control parameter for the actual companion deformation,
and then transport its complete numerical range through the periods.
All metrics and coefficient vectors below are the original ones.

Fix the original cyclic algebra, admitted degree and positive Gram
\[
 \begin{gathered}
 E=\mathbb C[S]/(\chi),\quad q=\deg\chi\ge1,\quad
 A=M_S,\quad G=G_N>0,\quad k\in\mathbb R,
 \\
 \quad e_0=[1],\quad
 \ell[P]=[S^{q-1}]\operatorname{rem}_{\chi}P.
 \end{gathered}
 \tag{LC1}
\]
Here $\chi$ is the original monic polynomial, $N\in\mathbb Z$ satisfies
$N\ge q-1$, and the deformation parameter is $t\in\mathbb C$.
Put \(R=e_0\ell\), \(A_t=A+tR\), and retain the actual control form
and relative, Laplacian and shifted endomorphisms
\[
 W_t=A_t^*G+GA_t-kG,\quad X_t=G^{-1}W_t,\quad
 L_t=A_t^2-kA_t,\quad B_t=A_t-\tfrac k2I.
 \tag{LC2}
\]
The control form has value $v^*W_tw=\langle v,X_tw\rangle_G$;
thus its exact form-to-operator map is $X_t=G^{-1}W_t$.
Here \(X_t\) is self-adjoint for
\(\langle v,w\rangle_G=v^*Gw\). Its finite operator norm is defined
by the supremum over the original nonzero vectors. Write
\(\epsilon_N=\|X_0\|_G\). This is an actual finite quantity, so
\(-\epsilon_NG\preceq W_0\preceq\epsilon_NG\) follows from the
spectral theorem for the displayed positive metric.

For the original reflection-stable source, its complete rank-two
calculation (TVB.22)--(TVB.30) evaluates this same quantity. In the
unchanged source-orthogonal columns \(b_j\) and norms \(\omega_j\),
put
\[
 a=b_N^*Gb_N,\quad d_+=b_{N+1}^*Gb_{N+1},\quad
 z_N=b_N^*Gb_{N+1},\quad D_N=\det K_N.
\]
For every integer $j\ge-1$ retain
$K_j=\sum_{r=0}^{j}b_rb_r^*/\omega_r$, with the empty sum
$K_{-1}=0$ on $E$, and put $D_j=\det K_j$.
The omitted-column kernel and determinant are
$\widetilde K_N=K_{N-1}+b_{N+1}b_{N+1}^*/\omega_{N+1}$ and
$\widetilde D_N=\det\widetilde K_N$.
At the first admitted degree $N=q-1$, the first $q-1$ monic
remainder columns are independent, so $K_{q-2}$ has rank $q-1$
and $D_{q-2}=0$. No inverse of this boundary kernel is used.
The retained identities give
\[
 \epsilon_N^2=\frac{ad_+-|z_N|^2}{\omega_N^2}
       =\alpha_{N+1}\frac{\Delta_N}{D_N},
 \quad\alpha_{N+1}=\frac{\omega_{N+1}}{\omega_N},
 \quad
 \Delta_N=D_{N+1}+D_{N-1}-D_N-\widetilde D_N.
 \tag{LC3}
\]
The omitted-column determinant \(\widetilde D_N\), every boundary
case and its positive source-minor sum are defined and proved in
that full chapter. To check the radius here directly, the original
control is
\(W_0=G(b_{N+1}b_N^*+b_Nb_{N+1}^*)G/\omega_N\).
Its relative self-adjoint operator has rank at most two and trace
zero; \(z_N\) is purely imaginary by the retained reflection phase.
Expansion of its squared trace gives
\(2(ad_+-|z_N|^2)/\omega_N^2\), so its two possible nonzero
eigenvalues are exactly \(\epsilon_N,-\epsilon_N\).
The two-column determinant identity proves the second equality
in (LC3). No estimate or limiting assumption is needed to evaluate it.

\subsection{The exact additional radius of the companion term}

The adjoint of \(R\) for the original metric is
\(R^{\dagger_G}=G^{-1}\ell^*e_0^*G\). Thus
\[
 X_t-X_0=tR+\bar tR^{\dagger_G}.
 \tag{LC4}
\]
For \(q>1\), \(\ell e_0=0\), hence \(R^2=0\). Put
\[
 f=G^{-1}\ell^*,\quad a_0=e_0^*Ge_0>0,\quad
 b_0=\ell G^{-1}\ell^*>0,\quad
 \sigma_N=\sqrt{a_0b_0}.
 \tag{LC5}
\]
The columns \(e_0,f\) are \(G\)-orthogonal because
\(e_0^*Gf=\overline{\ell e_0}=0\). They are independent and their
Gram is exactly \(\operatorname{diag}(a_0,b_0)\). In these original
unscaled columns the matrix of (LC4) is
\[
 \begin{pmatrix}0&t b_0\\\bar t a_0&0\end{pmatrix}.
 \tag{LC6}
\]
Indeed \(Rf=b_0e_0\), \(R e_0=0\),
\(R^{\dagger_G}e_0=a_0 f\) and
\(R^{\dagger_G}f=0\). On their \(G\)-orthogonal complement both
terms vanish. The characteristic polynomial on the displayed
two-dimensional space is \(x^2-|t|^2a_0b_0\). Therefore the exact
additional operator norm is \(|t|\sigma_N\); all remaining
eigenvalues are zero. This retains the two original vector lengths,
with no choice of unit vectors.

For \(q=1\), instead \(R=I_E\), so (LC4) equals
\(2\operatorname{Re}(t)I_E\) and its exact norm is
\(2|\operatorname{Re}t|\). Consequently the entirely specified
quantity
\[
 E_N(t)=\epsilon_N+
 \begin{cases}|t|\sigma_N,&q>1,\\
               2|\operatorname{Re}t|,&q=1
 \end{cases}
 \quad\text{satisfies}\quad
 \|X_t\|_G\le E_N(t),\qquad
 -E_N(t)G\preceq W_t\preceq E_N(t)G.
 \tag{LC7}
\]
The inequality is the triangle inequality for the same operator
norm. It is not an assertion that equality always holds. The two
summands in (LC7) are the actual norms just computed, and the
original \(\epsilon_N\) remains available through (LC3).
The full Laplacian change is also exact:
\[
 L_t-L_0=t(AR+RA-kR)+t^2R^2.
 \tag{LC8}
\]
Thus the quadratic term vanishes for \(q>1\), while for \(q=1\)
it is the retained scalar \(t^2I_E\).

\subsection{An explicit numerical-range enclosure with no missing input}

Expansion of (LC2), without a commutation assumption, proves
\[
 L_t^*G-GL_t=A_t^*W_t-W_tA_t,\qquad
 GL_t=W_tB_t-B_t^*GB_t-\tfrac{k^2}{4}G.
 \tag{LC9}
\]
For the second equality insert
\(W_t=B_t^*G+GB_t\), so the right side is
\(GB_t^2-k^2G/4=GL_t\). The first equality follows by expansion;
the two middle terms \(A_t^*GA_t\) cancel.

For each unchanged nonzero \(v\in E\), define the homogeneous
observations
\[
 z=\frac{v^*GL_tv}{v^*Gv},\quad
 r=\frac{\|B_tv\|_G}{\|v\|_G},\quad
 a+ib=\frac{\langle v,X_tB_tv\rangle_G}{\|v\|_G^2}.
 \tag{LC10}
\]
Cauchy--Schwarz and (LC7) give
\(a^2+b^2\le E_N(t)^2r^2\). The second identity in (LC9) gives
\(\operatorname{Re}z=a-r^2-k^2/4\) and
\(\operatorname{Im}z=b\). Write \(e=E_N(t)\). First,
\(a\le er\) implies
\(\operatorname{Re}z\le(e^2-k^2)/4-(r-e/2)^2\).
Second, direct multiplication gives
\[
 e^2\left(\frac{e^2-k^2}{4}-\operatorname{Re}z\right)
    -(\operatorname{Im}z)^2
   =(a-e^2/2)^2+e^2r^2-a^2-b^2\ge0.
\]
We have proved the actual enclosure
\[
 \boxed{\operatorname{Re}z\le\frac{E_N(t)^2-k^2}{4},\qquad
 (\operatorname{Im}z)^2\le E_N(t)^2
       \left(\frac{E_N(t)^2-k^2}{4}-\operatorname{Re}z\right).}
 \tag{LC11}
\]
It includes every eigenvalue by taking its nonzero eigenvector;
no diagonalizability is needed. It includes the case \(E_N(t)=0\).
Every quantity in (LC11) is supplied by the original finite metric,
the monic coefficient functional and the actual deformation parameter.

\subsection{The complete period transport}

For fixed \(u\ne0\), the full determinant proof makes the original
period matrix \(\Pi(u,t)\) invertible at every parameter. Put
\[
 \widetilde G=\Pi^{-*}G\Pi^{-1},\quad
 \widetilde A_t=\Pi A_t\Pi^{-1},\quad
 \widetilde L_t=\Pi L_t\Pi^{-1},\quad
 \widetilde W_t=\Pi^{-*}W_t\Pi^{-1}.
 \tag{LC12}
\]
Direct multiplication gives
\(\widetilde W_t=\widetilde A_t^*\widetilde G+
\widetilde G\widetilde A_t-k\widetilde G\) and
\(\widetilde L_t=\widetilde A_t^2-k\widetilde A_t\).
For \(x=\Pi v\), its norm and the complete numerical-range quotient
are exactly those in (LC10). Its control quotient and adjoint-defect
identity are likewise transported by the literal congruences in
(LC12). Thus (LC11) applies unchanged in that period realization.

The transported fixed action is \(\Pi A\Pi^{-1}\); the displayed
deformed action differs from it by \(t\Pi R\Pi^{-1}\), whose
control contribution was computed in (LC4)--(LC7). The arithmetic
source realizes precisely this change through
\(D_t^{(k)}=D^{(k)}+t r_NRj_E\), with
\(j_Er_N=I\) and the original boundary primitive retained, as
proved in (SC7)--(SC8). These equalities supply the source,
coefficient and period maps for the same Laplacian control.
To retain its complete primitive, let $\mathcal C_{\rm prim}$
be the original primitive term of the tensor complex and write
$D_0:\mathcal C_{\rm prim}\to\mathcal C_{\rm prim}$ for its
original Euler action of cochain degree zero. Thus
$K_N^{\rm prim}:E\to\mathcal C_{\rm prim}$ and
$D^{(k)}d=dD_0$ on this term. The original boundary observation
$j_Ed=0$ makes the added finite-rank action vanish there, so
$D_t^{(k)}d=dD_0$ as well. Put $K=K_N^{\rm prim}$.
Using $D_t^{(k)}r_N=r_NA_t+dK$, expand both terms of the square:
\[
\begin{split}
 (D_t^{(k)})^2r_N
 &=D_t^{(k)}(r_NA_t)+D_t^{(k)}dK\\
 &=(r_NA_t+dK)A_t+dD_0K
   =r_NA_t^2+d(KA_t+D_0K).
\end{split}
\]
Subtract $kD_t^{(k)}r_N=kr_NA_t+kdK$ with its original sign.
This proves the complete quadratic source comparison
\[
 \boxed{\bigl((D_t^{(k)})^2-kD_t^{(k)}\bigr)r_N-r_NL_t
       =d\bigl(D_0K_N^{\rm prim}+K_N^{\rm prim}A_t
                                      -kK_N^{\rm prim}\bigr).}
 \tag{LC12a}
\]
The term $K_N^{\rm prim}A_t=K_N^{\rm prim}A+
tK_N^{\rm prim}R$ retains the full deformation contribution.
All terms inside the displayed primitive have the same target
$\mathcal C_{\rm prim}$; the original differential, source action
and boundary class remain explicit.
No smallness uniform in packet size is claimed for the fully
calculated \(\epsilon_N\) or \(\sigma_N\).

\subsection{The exact finite comparison with constituent curvature}

Let $I:F\hookrightarrow E$ be an actual injective coefficient map
with $AI=IA_F$ and $0<p:=\dim F<q$. In particular $q>1$, so
the radius $\sigma_N$ in (LC5) applies. Retain the original
orthogonal projection and its two scalar observations:
\[
\begin{gathered}
 G_F=I^*GI,\qquad P_F=I G_F^{-1}I^*G,\qquad L=\ell I,\qquad
 f=G^{-1}\ell^*,\\
 a_\perp=e_0^*G(1-P_F)e_0,\qquad
 b_F=L G_F^{-1}L^*=\ell I G_F^{-1}I^*\ell^*.
\end{gathered}\tag{LC13}
\]
The injection and positive $G$ give $G_F>0$. Multiplication proves
$P_F^2=P_F$ and $P_F^*G=GP_F$, and its image is exactly $I(F)$.
Consequently $a_\perp=\|(1-P_F)e_0\|_G^2$ and
$b_F=\|P_Ff\|_G^2$, the second identity following from
$I^*Gf=I^*\ell^*$, with the same column $f$ as in (LC5).

Both numbers are strictly positive on this actual constituent.
Indeed $A$-invariance implies invariance under every polynomial in
$A$, so $I(F)$ is an ideal of $E$. Its inverse image in
$\mathbb C[S]$ is $(g)$ for a unique monic divisor $g$ of $\chi$.
Writing $\chi=gh$, the map $[a]_h\mapsto[ga]_\chi$ is an
isomorphism onto $I(F)$: it is onto by the ideal description, and
$gh\mid ga$ implies $h\mid a$ by cancellation in the polynomial
domain. Thus $\deg h=p$ and the complete original ideal has basis
$[g],[Sg],\ldots,[S^{p-1}g]$. The last vector is monic of degree
$q-1$, proving $L\ne0$; positivity of $G_F^{-1}$ gives $b_F>0$.
An ideal containing $e_0=[1]$ is all of $E$. Properness therefore
gives $e_0\notin I(F)$ and $a_\perp>0$. Every repeated factor
remains in this ideal calculation.

At the same $u\ne0$ and parameter $t$, let $m(t)$ and $M(t)$ be
the actual smallest and largest eigenvalues of
$B(t)=G^{-1}\Pi(t)^*\Pi(t)$. This operator is $G$-self-adjoint
and positive because $v^*GB(t)v=\|\Pi(t)v\|_{I_q}^2>0$ for
each nonzero $v$. Its finite spectrum therefore gives
$0<m(t)\le M(t)<\infty$ and the quadratic-form inequalities
$m(t)G\preceq\Pi(t)^*\Pi(t)\preceq M(t)G$.
These are evaluated finite comparison quantities, with no bound
uniform in the arithmetic inputs imposed on them.

Put $Y=\Pi I$, $H_F=Y^*Y$, and
$P^{\rm per}=YH_F^{-1}Y^*$, the orthogonal projection in the
specified period-coordinate metric $I_q$. Its residual is
$z=(1-P^{\rm per})\Pi e_0$. The exact comparison is
\[
 \boxed{m(t)a_\perp\le\|z(t)\|_{I_q}^2\le M(t)a_\perp,
 \qquad
 \frac{b_F}{M(t)}\le L H_F(t)^{-1}L^*
                         \le\frac{b_F}{m(t)}.}
 \tag{LC14}
\]
To prove the first pair, orthogonal projection gives
$\|z\|_{I_q}^2=\min_{x\in F}\|\Pi(e_0-Ix)\|_{I_q}^2$.
The same minimization in the original metric has value
$a_\perp$, attained at $x=G_F^{-1}I^*Ge_0$.
The lower quadratic-form bound holds for every $x$, so its minimum
is at least $m(t)a_\perp$. Evaluating the upper bound at the
displayed original minimizer gives the upper inequality.
Compression by $I$ gives $m(t)G_F\preceq H_F\preceq M(t)G_F$.
Inversion of positive forms, followed by the unchanged row $L$,
proves the second pair. Explicitly, for any positive matrix $T$
and column $w$, completing the square proves
$w^*T^{-1}w=\max_x(2\operatorname{Re}(w^*x)-x^*Tx)$.
Applying this formula to the three positive matrices verifies the
inverse order and both scalar factors without a commutation assumption.

The coefficient of the actual constituent curvature in (SC16) is
$c_F(t)=\|z(t)\|_{I_q}^2 L H_F(t)^{-1}L^*/|u|^2$.
For completeness, the period equation and $AI=IA_F$ give
$Y'=-YA_F/u-t\Pi e_0L/u$, hence
$(1-P^{\rm per})Y'=-tzL/u$.
Holomorphy of $Y$ gives $\partial_tH_F=Y^*Y'$ and
$\partial_{\bar t}H_F=Y'^*Y$; differentiating its inverse in
the determinant derivative proves
$\partial_{\bar t}\partial_t\log\det H_F
=\operatorname{Tr}(H_F^{-1}Y'^*(1-P^{\rm per})Y')$.
Substitution yields the full equality and the bounds
\[
\begin{gathered}
 \partial_{\bar t}\partial_t\log\det H_F(t)=|t|^2c_F(t),\\
 \boxed{\frac{m(t)}{M(t)}\frac{a_\perp b_F}{|u|^2}
       \le c_F(t)\le
       \frac{M(t)}{m(t)}\frac{a_\perp b_F}{|u|^2}.}
\end{gathered}\tag{LC15}
\]
Every factor in (LC14) is positive, so multiplying its appropriate
two lower or two upper bounds proves (LC15). In particular the
coefficient is positive even at $t=0$, while the actual curvature
at that parameter is zero through its retained factor $|t|^2$.

The comparison to the Laplacian radius has an exact two-term loss.
Orthogonality of $P_F$ gives
$a_0=\|P_Fe_0\|_G^2+a_\perp$ and
$b_0=b_F+\|(1-P_F)f\|_G^2$. Multiplying with every term
retained therefore proves
\[
 \boxed{\sigma_N^2-a_\perp b_F
  =\|P_Fe_0\|_G^2 b_0
      +a_\perp\|(1-P_F)G^{-1}\ell^*\|_G^2\ge0.}
 \tag{LC16}
\]
Combining (LC15)--(LC16) and the exact additional norm (LC4)--(LC6)
gives, on this same original constituent,
\[
 c_F(t)\le\frac{M(t)}{m(t)}\frac{\sigma_N^2}{|u|^2},\qquad
 \partial_{\bar t}\partial_t\log\det H_F(t)
 \le\frac{M(t)}{m(t)}
           \frac{\|X_t-X_0\|_G^2}{|u|^2}.
\]
Thus the companion control, the original constituent projections,
and its period curvature have an explicit finite comparison with
the complete loss and both metric comparison eigenvalues retained.
The proof assigns no uniform estimate to $M(t)/m(t)$ or $\sigma_N$.

\end{document}
